\begin{table}[t]
\centering
\begin{tabular}{lccc}
\hline
Comparison & J@32 & chance & ceiling \\
\hline
Llama base(16s) $\leftrightarrow$ FT(16s) & 0.524 & 0.031 & $\sim$0.71 \\
Llama FT: 0-shot $\leftrightarrow$ 16-shot (same weights) & 0.391 & 0.031 & $\sim$0.71 \\
Llama FT $\leftrightarrow$ TS taught (cross-model) & 0.164\textsuperscript{a} & 0.065 & -- \\
Llama FT $\leftrightarrow$ FT+format-install & 0.684 & 0.031 & $\sim$0.71 \\
TS1B-latent $\leftrightarrow$ elicited child (op) & 0.455--0.488 & 0.031 & 0.684 \\
TS1B-latent $\leftrightarrow$ elicited child (bridge) & 0.524 & 0.031 & 0.684 \\
Elicited-1M $\leftrightarrow$ taught-1M (same base) & 0.231\textsuperscript{b} & 0.031 & 0.684/0.561 \\
\hline
\end{tabular}
\caption{Circuit-membership overlap (attribution top-32 Jaccard). \textsuperscript{a}J@64, score Spearman $-0.43$. \textsuperscript{b}score Spearman $-0.11$: same substrate, same task, different mechanism. Fine-tuning a latent-capable model moves its circuit less than changing the prompt regime does (0.455--0.524 vs 0.391); teaching builds elsewhere. All maps performing-regime-guarded; decisions.md 2026-08-24..2026-09-07.}
\label{tab:circuit_overlap}
\end{table}
